
\begin{lemma}\label{MAT-L-PEO of complete graph}
Let $ G = K_{\ell} $ be a complete graph and $ W \subseteq V_{G} $. 
\begin{enumerate}[(1)]
\item\label{MAT-L-PEO of complete graph 1} 
Let $ \lambda $ be an MAT-labeling of $ G$. If $ (v_{1}, \dots, v_{r}) $ is an MAT-PEO  of $ (G[W], \lambda|_{E_{G[W]}}) $, then $ (v_{1}, \dots, v_{r}) $ can be extended to an MAT-PEO $ (v_{1}, \dots, v_{r}, \dots, v_{\ell}) $ of $ (G,\lambda) $. 
\item\label{MAT-L-PEO of complete graph 2} 

If $ \lambda_{W} $ is an MAT-labeling of $ G[W] $, then $ \lambda_{W} $ can be extended to an MAT-labeling of $ G $. 

\end{enumerate}
As a consequence, a complete graph always has an MAT-labeling (and an MAT-PEO) which can be constructed inductively from any vertex of the graph.
\end{lemma}
\begin{proof}
\ref{MAT-L-PEO of complete graph 1}
We proceed by induction on $ \ell = |V_{G}| $. 
When $ \ell = 1 $, we have nothing to prove. 
Now suppose $ \ell \geq 2 $. 
If $ W = V_{G} $, the assertion holds trivially. 
Suppose $ W \subsetneq V_{G} $. 
Let $ e_{0}\in E_G $ be the edge with maximal label. 
Since $ \max\Set{\lambda(e) \in \mathbb{Z}_{>0} | e \in E_{G[W]}} < \lambda(e_{0}) = \ell-1 $, at least one endvertex of $ e_{0} $, say $ v_{\ell} $  does not belong to $ W $. 
By Lemma~\ref{lem:MATS-existence}\ref{MATS-existence 1}, $ v_{\ell} $ is MAT-simplicial in $ G $. 
By the induction hypothesis, there exists an MAT-PEO $ (v_{1}, \dots, v_{r}, \dots, v_{\ell-1}) $ of $ G\setminus v_{\ell} $. 
Hence $ (v_{1}, \dots, v_{\ell}) $ is a desired ordering. 

\ref{MAT-L-PEO of complete graph 2} 
Without loss of generality, we may assume $ \ell \geq 2 $ and $|W| = \ell - 1 $. 
By Theorem~\ref{MAT-labeling ordering}, there exists an MAT-PEO $ (v_{1}, \dots, v_{\ell-1}) $ of $ (G[W], \lambda_{W}) $.
Let $ v_{\ell} $ denote the vertex in $ V_{G}\setminus W $. 
We define a labeling $ \lambda $ of $ G $ by 
\begin{align*}
\lambda(e) \coloneqq
\begin{cases}
\lambda_{W}(e) & \text{ if } e \in E_{G[W]}; \\
i &  \text{ if } e = \{v_{i}, v_{\ell}\} \text{ for } i \in  [\ell -1]. 
\end{cases}
\end{align*}

We will show that $ v_{\ell} $ is MAT-simplicial in $ (G,\lambda) $. 
Firstly, \ref{definition MAT-simplicial 1} is clear since $ G $ is complete. 
Secondly, we have $ \Set{\lambda(\{v_{i}, v_{\ell}\}) | i \in [\ell-1]} = [\ell-1] $ and hence \ref{definition MAT-simplicial 2} holds. 
Thirdly, for $ 1\le i < j < \ell $, we have $ \lambda(\{v_{i}, v_{j}\}) \leq j-1 < j = \lambda(\{v_{j}, v_{\ell}\}) $, which shows \ref{definition MAT-simplicial 3}. 
Therefore $ v_{\ell} $ is MAT-simplicial in $ (G,\lambda) $. 
Thus $ (v_{1}, \dots, v_{\ell} )$ is an MAT-PEO in $ (G, \lambda) $ and $ \lambda $ is an MAT-labeling by Proposition \ref{MAT-simplicial}.
\end{proof}


\begin{lemma}\label{merge}
Let $ G = K_{\ell} $. 
Suppose that $ V_{G} = A \cup B $ and there exist MAT-labelings $ \lambda_{A}, \lambda_{B}, \lambda_{A \cap B} $ of $ G[A], G[B], G[A\cap B] $, respectively such that $ \lambda_{A}|_{E_{G[A\cap B]}} = \lambda_{B}|_{E_{G[A\cap B]}} = \lambda_{A \cap B} $. 
Then there exists an MAT-labeling $ \lambda $ of $ G $ such that $ \lambda|_{E_{G[A]}} = \lambda_{A} $ and $ \lambda|_{E_{G[B]}} = \lambda_{B} $. 
\end{lemma}
\begin{proof}
By Lemma \ref{MAT-L-PEO of complete graph}\ref{MAT-L-PEO of complete graph 1}, there exists an MAT-PEO $ (a_{1}, \dots, a_{p}) $ of $ G[A \cap B] $ and its extensions $ (a_{1}, \dots, a_{p}, a_{p+1}, \dots, a_{p+q}) $ of $ G[A] $ and $ (a_{1}, \dots, a_{p}, b_{1}, \dots, b_{r}) $ of $ G[B] $ (where $p+q+r=\ell$). 
Define a labeling $ \lambda \colon E_{G} \to \mathbb{Z}_{>0} $ by 
\begin{align*}
\lambda(e) \coloneqq \begin{cases}
\lambda_{A}(e) & \text{ if } e \in E_{G[A]}; \\
\lambda_{B}(e) & \text{ if } e \in E_{G[B]}; \\
p + i + j - 1 & \text{ if } e = \{a_{p+i}, b_{j}\} \quad (i \in [q], \,  j \in[  r]). 
\end{cases}
\end{align*}

We claim that $ \lambda $ is a desired MAT-labeling of $ G $ by induction on $ r $. 
If $ r = 0 $, then $ \lambda = \lambda_{A} $ and hence the claim holds. 
Suppose $ r \geq 1 $. 
We will prove that $ b_{r} $ is MAT-simplicial in $ (G, \lambda) $. 
The condition \ref{definition MAT-simplicial 1} is clear. 
Since $ b_{r} $ is MAT-simplicial in~$ (G[B], \lambda_{B}) $, 
\begin{align*}
\Set{\lambda(\{a_{i}, b_{r}\}) | i \in [p ] } \cup \Set{\lambda(\{b_{j}, b_{r}\}) | j \in[  r-1]} =[ p+r-1]. 
\end{align*}
By the definition of $ \lambda $ we have $ \lambda(\{a_{p+i}, b_{r}\}) = p+i+r-1 \ (i \in [q]) $. 
Therefore~$ \Set{\lambda(\{v,b_{r}\}) | v \in N_{G}(b_{r})} = [\ell - 1] $ and hence \ref{definition MAT-simplicial 2} holds. 

Next we show \ref{definition MAT-simplicial 3}, \ie  $ \lambda(\{u,v\}) < \max\{\lambda(\{u,b_{r}\}), \lambda(\{v,b_{r}\})\} $ for any distinct vertices $ u,v \in N_{G}(b_{r}) $. 
It is clear when $ u,v \in B $ since $ b_{r} $ is MAT-simplicial in~$ (G[B], \lambda_{B}) $. 
Consider the case $ u=a_{p+i} \in A\setminus B, v=a_{j} \in A \cap B$ for some $ i,j $ with $ p+i > j $. 
Then $ \lambda(\{a_{p+i}, a_{j}\}) \leq p+i-1 < p+i+r-1 = \lambda(\{a_{p+i}, b_{r}\}) $ since $ a_{p+i} $ is MAT-simplicial in $ (G, \lambda|_{E_{G[\{a_{1}, \dots, a_{p+i}\}]}}) $. 
Now consider $ u=a_{p+i} \in A\setminus B, v = b_{j} \in B\setminus A $ for some $ i,j $ with $ 1 \leq i \leq q $ and $ 1 \leq j < r $. 
Then $ \lambda(\{a_{p+i}, b_{j}\}) = p+i+j-1 < p+i+r-1 = \lambda(\{a_{p+i}, b_{r}\}) $. 
Thus \ref{definition MAT-simplicial 3} holds and $ b_{r} $ is an MAT-simplicial vertex of $ (G, \lambda) $. 

By the induction hypothesis, $ \lambda|_{E_{G\setminus b_{r}}} $ is an MAT-labeling. 
Using Proposition \ref{MAT-simplicial}, we conclude that $ \lambda $ is an MAT-labeling. 
\end{proof}
